\documentclass[aspectratio=169,12pt]{beamer}
% Estilo canónico para las presentaciones de Métodos Numéricos 2026-2.
\usepackage[T1]{fontenc}
\usepackage[utf8]{inputenc}
\usepackage[spanish,es-nodecimaldot]{babel}
\usepackage{lmodern}
\usepackage{amsmath,amssymb,mathtools}
\usepackage{booktabs}
\usepackage{graphicx}
\usepackage{xcolor}
\usepackage{hyperref}

\definecolor{upazul}{HTML}{003B70}
\definecolor{upazulclaro}{HTML}{EAF2F8}
\definecolor{updorado}{HTML}{C5962E}
\definecolor{upgris}{HTML}{4A5560}
\definecolor{upverde}{HTML}{167C5A}
\definecolor{uprojo}{HTML}{B33A3A}

\setbeamercolor{normal text}{fg=black,bg=white}
\setbeamercolor{structure}{fg=upazul}
\setbeamercolor{frametitle}{fg=white,bg=upazul}
\setbeamercolor{title}{fg=upazul}
\setbeamercolor{subtitle}{fg=upgris}
\setbeamercolor{block title}{fg=white,bg=upazul}
\setbeamercolor{block body}{fg=black,bg=upazulclaro}
\setbeamercolor{alerted text}{fg=uprojo}
\setbeamercolor{example text}{fg=upverde}

\setbeamerfont{title}{size=\Huge,series=\bfseries}
\setbeamerfont{subtitle}{size=\Large}
\setbeamerfont{frametitle}{size=\LARGE,series=\bfseries}
\setbeamerfont{normal text}{size=\large}
\setbeamerfont{block title}{series=\bfseries}

\setbeamertemplate{navigation symbols}{}
\setbeamertemplate{itemize item}{\color{updorado}$\blacktriangleright$}
\setbeamertemplate{itemize subitem}{\color{upazul}$\bullet$}
\setbeamertemplate{blocks}[rounded][shadow=false]

\setbeamertemplate{frametitle}{%
  \nointerlineskip
  \begin{beamercolorbox}[wd=\paperwidth,ht=0.78cm,dp=0.24cm,leftskip=0.55cm]{frametitle}
    \usebeamerfont{frametitle}\insertframetitle
  \end{beamercolorbox}}

\setbeamertemplate{footline}{%
  \leavevmode
  \hbox{%
    \begin{beamercolorbox}[wd=.82\paperwidth,ht=2.7ex,dp=1.2ex,leftskip=.5cm]{author in head/foot}
      \color{upgris}\scriptsize Métodos Numéricos 2026-2 \hspace{0.8em} | \hspace{0.8em} Universidad del Pacífico
    \end{beamercolorbox}%
    \begin{beamercolorbox}[wd=.18\paperwidth,ht=2.7ex,dp=1.2ex,rightskip=.5cm plus1fil]{date in head/foot}
      \color{upazul}\scriptsize\hfill\insertframenumber/\inserttotalframenumber
    \end{beamercolorbox}}
  \vskip0pt}

\newcommand{\concepto}[1]{\textcolor{upazul}{\textbf{#1}}}
\newcommand{\resultado}[1]{\textcolor{upverde}{\textbf{#1}}}
\newcommand{\advertencia}[1]{\textcolor{uprojo}{\textbf{#1}}}
\newcommand{\codigo}[1]{\colorbox{upazulclaro}{\strut\texttt{#1}}}
\newcommand{\fuente}[1]{\vfill\hfill{\tiny\color{upgris}#1}}


\title{Mínimos cuadrados}
\subtitle{QR, condicionamiento y aplicaciones}
\author{Oswaldo José Velásquez Castañón}
\institute{Universidad del Pacífico}
\date{Métodos Numéricos -- 2026-2}
\newcommand{\R}{\mathbb{R}}
\newcommand{\mat}[1]{\begin{bmatrix}#1\end{bmatrix}}
\newcommand{\norm}[1]{\left\lVert#1\right\rVert}
\begin{document}
\begin{frame}[plain]
\begin{columns}[T,onlytextwidth]
\begin{column}{.70\textwidth}
\vspace{.7cm}{\usebeamerfont{title}\color{upazul}\inserttitle\par}
\vspace{.35cm}{\usebeamerfont{subtitle}\color{upgris}\insertsubtitle\par}
\vspace{1cm}{\large\insertauthor\par}{\normalsize\insertdate\par}
\end{column}
\begin{column}{.25\textwidth}\vspace{.55cm}\includegraphics[width=\linewidth]{UPlogo.pdf}\end{column}
\end{columns}
\end{frame}

\begin{frame}{Objetivos de la clase}
\begin{itemize}
\item Resolver mínimos cuadrados mediante factorización QR.
\item Explicar la pérdida de precisión al formar $A^TA$.
\item Reconocer columnas casi dependientes y problemas de escala.
\item Formular ajustes ponderados y regularizados.
\item Distinguir un ajuste en niveles de uno en logaritmos.
\end{itemize}
Usaremos la misma recta y los mismos datos de la clase anterior.
\end{frame}

\begin{frame}{Ecuaciones normales y cómputo}
La caracterización matemática es
\[
A^TA\theta=A^Tb.
\]
Si $A$ tiene rango columna completo, podemos formar $A^TA$, aplicar Cholesky
y resolver dos sistemas triangulares.
\begin{block}{Dificultad numérica}
El producto $A^TA$ puede perder información cuando las columnas de $A$ son
casi dependientes. QR trabaja directamente con $A$.
\end{block}
La equivalencia algebraica entre dos métodos no garantiza la misma precisión.
\end{frame}

\begin{frame}{Condicionamiento}
Para una matriz rectangular de rango columna completo,
\[
\kappa_2(A)=\frac{\sigma_{\max}(A)}{\sigma_{\min}(A)}.
\]
Los valores propios de $A^TA$ son los cuadrados de los valores singulares de $A$.
Así,
\[
\boxed{\kappa_2(A^TA)=\kappa_2(A)^2.}
\]
Si $\kappa_2(A)=10^6$, el sistema normal tiene condición $10^{12}$.
Este ejemplo describe el condicionamiento, no una cifra exacta de error final.
\end{frame}

\begin{frame}{Columnas casi dependientes}
Considere
\[
A_\varepsilon=\mat{1&1\\0&\varepsilon},\qquad
A_\varepsilon^TA_\varepsilon
=\mat{1&1\\1&1+\varepsilon^2}.
\]
La independencia de las columnas depende de $\varepsilon$.
Cuando $\varepsilon^2$ queda por debajo de la precisión relativa de la máquina,
el cálculo de $1+\varepsilon^2$ puede redondearse a $1$.
\begin{block}{Consecuencia}
Una matriz de rango completo puede producir un sistema normal singular
en aritmética de máquina.
\end{block}
\end{frame}

\begin{frame}{Factorización QR completa}
Sea $A\in\R^{m\times n}$, $m\geq n$, de rango $n$. Escribimos
\[
A=Q\mat{R\\0},\qquad
Q^TQ=I_m,\qquad R\in\R^{n\times n}.
\]
La matriz $R$ es triangular superior y no singular. Separamos
\[
Q^Tb=\mat{c\\d},\qquad c\in\R^n,\quad d\in\R^{m-n}.
\]
Householder y Givens permiten construir esta factorización con transformaciones ortogonales.
\end{frame}

\begin{frame}{Reducción a un sistema triangular}
Como las matrices ortogonales preservan la norma,
\[
\begin{aligned}
\norm{A\theta-b}_2^2
&=\norm{Q^T(A\theta-b)}_2^2\\
&=\norm{R\theta-c}_2^2+\norm{d}_2^2.
\end{aligned}
\]
El segundo término no depende de $\theta$. Por tanto,
\[
\boxed{R\theta^*=c},\qquad
\boxed{\min_\theta\norm{A\theta-b}_2=\norm{d}_2}.
\]
Una sustitución hacia atrás recupera los coeficientes.
\end{frame}

\begin{frame}{QR reducida}
Si $A=Q_1R$ con $Q_1\in\R^{m\times n}$ y $Q_1^TQ_1=I_n$,
\[
R\theta^*=Q_1^Tb.
\]
La proyección y el residuo son
\[
\widehat b=Q_1Q_1^Tb,\qquad r=b-\widehat b.
\]
Para el ajuste basta almacenar las primeras $n$ columnas de $Q$ o aplicar
sus reflectores directamente a $b$.
\begin{block}{Conexión con el capítulo de QR}
La factorización es la misma. Ahora la usamos para resolver un problema de aproximación.
\end{block}
\end{frame}

\begin{frame}{Ejemplo: QR de la matriz de diseño}
Para $A=\mat{1&0\\1&1\\1&2\\1&3}$, una QR reducida es
\[
Q_1=\mat{
\frac12&-\frac{3}{2\sqrt5}\\
\frac12&-\frac{1}{2\sqrt5}\\
\frac12& \frac{1}{2\sqrt5}\\
\frac12& \frac{3}{2\sqrt5}},
\qquad R=\mat{2&3\\0&\sqrt5}.
\]
Se verifica $Q_1^TQ_1=I_2$ y $Q_1R=A$.
La segunda columna ortogonaliza el predictor respecto de la constante.
\end{frame}

\begin{frame}{Ejemplo: sustitución hacia atrás}
Con $b=(1,2,2,4)^T$,
\[
Q_1^Tb=\mat{4.5\\4.5/\sqrt5}.
\]
Resolvemos
\[
\mat{2&3\\0&\sqrt5}\mat{c\\a}
=\mat{4.5\\4.5/\sqrt5}.
\]
Primero $a=4.5/5=0.9$. Después,
\[
c=\frac{4.5-3(0.9)}2=0.9.
\]
Obtenemos la misma recta $\widehat y=0.9+0.9x$.
\end{frame}

\begin{frame}{Algoritmo de mínimos cuadrados}
Entrada: matriz $A$ de rango columna completo y vector $b$.
\begin{enumerate}
\item Factorizar $A$ mediante Householder o Givens.
\item Aplicar las transformaciones ortogonales a $b$.
\item Resolver $R\theta=c$ por sustitución hacia atrás.
\item Calcular $r=b-A\theta$ en las variables originales.
\item Revisar la norma del residuo y $A^Tr$.
\end{enumerate}
No se necesita formar una inversa explícita.
\end{frame}

\begin{frame}{Centrado del predictor}
Tomamos $t_i=x_i-\bar x$ y escribimos
\[
\widehat y=\beta_0+\beta_1t.
\]
Como $\sum_i t_i=0$, las columnas $1$ y $t$ son ortogonales:
\[
A_t^TA_t=\mat{m&0\\0&\sum_i t_i^2}.
\]
En nuestro ejemplo, $\bar x=1.5$, $\beta_0=2.25$ y $\beta_1=0.9$.
Se recupera el intercepto original mediante
\[
c=\beta_0-\beta_1\bar x=0.9.
\]
\end{frame}

\begin{frame}{Escala de las columnas}
Si los predictores tienen unidades o magnitudes muy diferentes, podemos usar
\[
t_i=\frac{x_i-\bar x}{s_x},\qquad s_x>0.
\]
\begin{itemize}
\item El escalado mejora el equilibrio entre las magnitudes de las columnas.
\item Debemos transformar los coeficientes al interpretar el modelo original.
\item El escalado no elimina una dependencia lineal verdadera.
\end{itemize}
En bases polinómicas, centrar y escalar resulta especialmente útil.
\end{frame}

\begin{frame}{Rango deficiente y SVD}
Si $A=U\Sigma V^T$, una solución de mínima norma es
\[
\theta^*=A^+b
=\sum_{j:\sigma_j>0}\frac{u_j^Tb}{\sigma_j}v_j.
\]
$A^+$ es la pseudoinversa. Los valores singulares pequeños amplifican las
perturbaciones de $b$.
\begin{block}{Rango numérico}
En cómputo se usa una tolerancia para decidir qué valores singulares son
significativos. QR con pivoteo también ayuda a detectar dependencia.
\end{block}
\end{frame}

\begin{frame}{Mínimos cuadrados ponderados}
Con pesos $w_i>0$, minimizamos
\[
\sum_i w_i(y_i-\widehat y_i)^2
=\norm{W^{1/2}(A\theta-b)}_2^2,
\quad W=\operatorname{diag}(w_1,\ldots,w_m).
\]
La formulación equivalente es
\[
\widetilde A=W^{1/2}A,\qquad \widetilde b=W^{1/2}b.
\]
Aplicamos QR a $\widetilde A$.
Si las observaciones tienen varianzas conocidas y errores independientes,
una elección habitual es $w_i=1/\sigma_i^2$.
\end{frame}

\begin{frame}{Regularización de Tikhonov}
Para limitar coeficientes grandes, consideramos
\[
J_\lambda(\theta)=\norm{A\theta-b}_2^2+\lambda\norm{\theta}_2^2,
\qquad\lambda>0.
\]
Las condiciones de optimalidad son
\[
(A^TA+\lambda I)\theta^*=A^Tb.
\]
Como $A^TA+\lambda I$ es definida positiva, el minimizador es único,
incluso si $A$ tiene rango deficiente.
El parámetro $\lambda$ equilibra ajuste y tamaño de los coeficientes.
\end{frame}

\begin{frame}{Regularización mediante QR}
La misma función objetivo puede escribirse como
\[
J_\lambda(\theta)
=\left\lVert
\mat{A\\\sqrt\lambda I}\theta-\mat{b\\0}
\right\rVert_2^2.
\]
Así resolvemos un problema aumentado de mínimos cuadrados.
\begin{block}{Penalización selectiva}
Si queremos penalizar solo algunos coeficientes, reemplazamos $I$ por una
matriz $L$ y usamos $\lambda\norm{L\theta}_2^2$.
La unicidad exige $\mathcal N(A)\cap\mathcal N(L)=\{0\}$.
\end{block}
\end{frame}

\begin{frame}{Demostración del mínimo regularizado}
Si $(A^TA+\lambda I)\theta^*=A^Tb$, entonces
\[
J_\lambda(\theta^*+h)-J_\lambda(\theta^*)
=\norm{Ah}_2^2+\lambda\norm{h}_2^2.
\]
El miembro derecho es positivo para $h\neq0$.
\begin{block}{Conclusión}
$\theta^*$ es el único mínimo global. La penalización añade curvatura positiva
en todas las direcciones.
\end{block}
Esta identidad permite verificar la solución sin invocar únicamente el gradiente.
\end{frame}

\begin{frame}{Modelo de potencia}
Un modelo $y=bx^a$, con $x,y>0$ y $b>0$, se transforma en
\[
\ln y=\ln b+a\ln x.
\]
Con $z_i=\ln x_i$ y $v_i=\ln y_i$, ajustamos
\[
v_i\approx\theta_0+\theta_1z_i,
\qquad b=e^{\theta_0},\quad a=\theta_1.
\]
\begin{block}{Objetivo que cambia}
El ajuste minimiza errores cuadrados en logaritmos.
Minimizar errores cuadrados en niveles define un problema distinto.
\end{block}
\end{frame}

\begin{frame}{Bases trigonométricas}
Para aproximar datos periódicos, podemos elegir
\[
\widehat y(x)=\theta_0+\theta_1\sin x+\theta_2\cos x.
\]
Entonces
\[
A=\mat{1&\sin x_1&\cos x_1\\\vdots&\vdots&\vdots\\1&\sin x_m&\cos x_m}.
\]
La selección de la base incorpora estructura al modelo.
Los coeficientes se obtienen con el mismo procedimiento QR.
\end{frame}

\begin{frame}{Validación del ajuste}
\begin{itemize}
\item Residuo: calcular $\norm{b-A\theta}_2$ en las unidades originales.
\item Optimalidad: comprobar que $A^T(b-A\theta)$ es pequeño respecto de las escalas.
\item Rango: examinar valores singulares o la factorización con pivoteo.
\item Modelo: buscar patrones sistemáticos en los residuos.
\item Predicción: evaluar con datos que no participaron en el ajuste.
\end{itemize}
La estabilidad del algoritmo y la adecuación del modelo son propiedades distintas.
\end{frame}

\begin{frame}{Actividad de clase}
Use $t_i=x_i-1.5$ para los datos de la clase y ajuste
$\beta_0+\beta_1t_i$.
\begin{enumerate}
\item Compruebe que las columnas $1$ y $t$ son ortogonales.
\item Recupere $\beta_0=2.25$ y $\beta_1=0.9$ mediante QR.
\item Penalice solo la pendiente con $\lambda=1$.
\item Verifique que la nueva pendiente es $0.75$.
\item Calcule el nuevo error de ajuste y explique el compromiso introducido.
\end{enumerate}
\end{frame}

\begin{frame}{Ejercicio original: ajuste trigonométrico}
Ajuste los datos mediante una combinación lineal de $1$, $\sin x$ y $\cos x$:
\begin{center}\begin{tabular}{rrrrrrr}\toprule
$x_i$ & 1.1 & 1.9 & 2.7 & 4.3 & 5.1 & 7.8\\
$y_i$ & 11.2 & 16.8 & 18.7 & 8.1 & 2.7 & 14.3\\\bottomrule
\end{tabular}\end{center}
\begin{enumerate}
\item Escriba la matriz de diseño y las ecuaciones normales.
\item Resuelva el problema usando QR.
\item Calcule los valores ajustados y la norma del residuo.
\end{enumerate}
\fuente{Notas del curso, capítulo Aproximación lineal, ejercicio 1}
\end{frame}

\begin{frame}{Ejercicio original: regularización}
Sea $A\in\R^{m\times n}$, $b\in\R^m$ y $\alpha>0$. Defina
\[
F(\theta)=\norm{A\theta-b}_2^2+\alpha\norm{\theta}_2^2.
\]
Demuestre que $F$ alcanza su único mínimo en la solución de
\[
(A^TA+\alpha I)\theta=A^Tb.
\]
Para esa solución, pruebe la identidad
\[
F(\theta+h)=F(\theta)+\norm{Ah}_2^2+\alpha\norm{h}_2^2.
\]
\fuente{Notas del curso, ejercicio 2, identidad algebraica corregida}
\end{frame}

\begin{frame}{Ejercicio original: sistema transpuesto}
Encuentre la solución en mínimos cuadrados de
\[
\begin{bmatrix}u&v\end{bmatrix}
\mat{3&2&1\\2&3&2}
=\begin{bmatrix}3&0&1\end{bmatrix}.
\]
\begin{enumerate}
\item Transponga la expresión y formule $A\theta\approx b$.
\item Obtenga y resuelva las ecuaciones normales.
\item Verifique que el residuo es ortogonal a las columnas de $A$.
\end{enumerate}
\fuente{Notas del curso, capítulo Aproximación lineal, ejercicio 3}
\end{frame}

\begin{frame}{Ejercicio original: modelo de potencia}
Las notas recogen datos de peso $W$ (gramos) y consumo de oxígeno $R$
(mililitros por hora) para larvas de \emph{Pachysphinx modesta}.
\begin{center}\small\begin{tabular}{rrrrrrrr}\toprule
$W$ & .017 & .087 & .085 & .020 & .233 & 4.83 & 5.53\\
$R$ & .154 & .296 & .260 & .180 & .537 & 4.66 & 6.94\\\bottomrule
\end{tabular}\end{center}
Ajuste el modelo $R=bW^a$ mediante logaritmos.
\begin{enumerate}
\item Obtenga $a$ y $b$ por mínimos cuadrados.
\item Calcule el error en logaritmos y el error en la escala original.
\end{enumerate}
\fuente{Datos reproducidos de las notas del curso, ejercicio 4, atribuidos allí a L. Schroeder}
\end{frame}

\begin{frame}{Ejercicio original: modelo ampliado}
Con los mismos datos, ajuste
\[
R=bW^aW^{c\ln W}.
\]
Tomar logaritmos da
\[
\ln R=\ln b+a\ln W+c(\ln W)^2.
\]
\begin{enumerate}
\item Construya la matriz de diseño de tres columnas.
\item Calcule los nuevos parámetros mediante QR.
\item Compare el error en logaritmos con el modelo anterior.
\item Explique qué garantiza agregar una columna al espacio de ajuste.
\end{enumerate}
\fuente{Notas del curso, capítulo Aproximación lineal, ejercicio 4(b)}
\end{frame}

\begin{frame}{Cierre de la clase}
\begin{itemize}
\item QR reduce mínimos cuadrados a un sistema triangular.
\item Formar $A^TA$ eleva al cuadrado el número de condición.
\item Centrado, escala y rango influyen en la calidad numérica.
\item Los pesos y la regularización modifican el objetivo de aproximación.
\end{itemize}
\vspace{.4cm}
El siguiente capítulo estudia interpolación y la construcción de funciones
que cumplen condiciones exactas en los nodos.
\end{frame}
\end{document}
